% Section content template for individual Educative lessons
% This file contains the actual content for a specific section
% Generated from Educative API section content

% Section: Summary: Designing a Meeting Scheduler
% Section ID: 5490132386709504
% Section Slug: summary-designing-a-meeting-scheduler
% Generated: 2025-12-12T13:06:48.691121

Now that you've completed the meeting scheduler case study, let's take a moment to reflect on and consolidate what we've learned. We 'll revisit the key system requirements, identify the core classes along with their responsibilities and relationships, and highlight the major design principles applied. We
'll also examine how objects interact within the system and walk through the overall workflow to understand how the components come together to achieve the desired functionality.

\subsection{Key requirements}\label{Esc7TnxY3ARMPBTaK2weO}
\\
The following are the primary functional and operational requirements for the meeting scheduler system:

\begin{enumerate}
\item
\protect\phantomsection\label{8tpq12r4RtkWyzKvfmz48}
Support a configurable number of meeting rooms with editable capacities.
\item
\protect\phantomsection\label{FFGlwFWlsUSCO4BRkYR8v}
Prevent room overbooking by checking availability and overlapping intervals.
\item
\protect\phantomsection\label{X_yfvyZRUYMS2mXhGCwW8}
Automatically notify participants when meetings are scheduled, updated, or canceled.
\item
\protect\phantomsection\label{BxVWsHrSJ3YHDm--YvCO9}
Track each participant's invitation response (accepted, declined, pending).
\item
\protect\phantomsection\label{Zbg7NhPtMEG16-CR5MU5L}
Provide each user with a personal calendar for meeting management.
\item
\protect\phantomsection\label{8tJX5Ri3G2CDDPhNdznqm}
Allow organizers to add or remove participants and notify them accordingly.
\item
\protect\phantomsection\label{l3KYYSTiBgXmRqihTWLTy}
Handle scheduling conflicts and overlapping meetings with clear feedback.
\item
\protect\phantomsection\label{B0XnuUnV7qIgTihY_5gIO}
Allow editing of meeting details with real-time updates and notifications.
\item
\protect\phantomsection\label{LoIxK-T-bXFUfFDrHtViN}
Allow organizers or authorized users to cancel meetings and inform attendees.
\item
\protect\phantomsection\label{oVY1Y_NqhTmy8pR8xTuQ7}
Remove a meeting from a participant's calendar if they decline or withdraw.
\item
\protect\phantomsection\label{w65QiVxfKI-dKBWmRG8r2}
Autoupdate calendars and send notifications when any meeting property changes.
\end{enumerate}

\subsection{System actors}\label{Nv3ea0K348OIzNj7aBQUQ}
\\
The system involves both primary and secondary actors who interact with or support the core functionalities:

\begin{itemize}
\item
\protect\phantomsection\label{SV3WAYGhADAqHfAU6PakD}
\textbf{Organizer:} Schedules, edits, or cancels meetings and manages the list of participants.
\item
\protect\phantomsection\label{-uad9sv99mQA2ynWqmZvS}
\textbf{Participant:} Receives invitations, accepts/declines meetings, and maintains their calendar.
\end{itemize}

\subsection{Important classes and relationships}\label{MwHfrjvXR7RgppNieLRdq}
\\
The system comprises several core classes, each with defined responsibilities and relationships.

\begin{center}
\begin{tabular}{|p{0.20\linewidth}|p{0.38\linewidth}|p{0.38\linewidth}|}
\hline
\textbf{Class} & \textbf{Description} & \textbf{Important Relationships} \\
\hline

User &
Represents a person who organizes or participates in meetings. &
Composes one Calendar; associated with multiple Meeting instances. \\
\hline

Interval &
Defines a time slot using start and end times. &
Aggregated in Meeting and MeetingRoom. \\
\hline

MeetingRoom &
Represents a meeting room with capacity and booked intervals. &
Aggregates Interval; associated with multiple Meeting instances. \\
\hline

Meeting &
Contains meeting details: time, room, participants, and RSVP statuses. &
Associated with User, MeetingRoom, and Interval; composed in Calendar. \\
\hline

Calendar &
Maintains a list of a user’s meetings. &
Composed in User; composed of multiple Meeting instances. \\
\hline

MeetingScheduler &
Central scheduler managing all meetings and rooms. &
Aggregates Meeting and MeetingRoom; handles booking, releasing, notifications. \\
\hline

Notification &
Sends invites and cancellations to users. &
Associated with User. \\
\hline

RSVPStatus &
Enum representing a participant’s response status. &
Used in Meeting for tracking participant responses. \\
\hline

\end{tabular}
\end{center}

\subsection{Design highlights}\label{NTm7DhN9WP_ZF4KiRASj1}
\\
This section summarizes the overall design approach and highlights how key principles and patterns were applied to develop the meeting scheduler system.

\subsubsection{Design approach}\label{nmjhY__HQlzSGFBzBJ_oJ}
\\
The system was designed using a bottom-up approach, starting with foundational classes like \texttt{User}, \texttt{MeetingRoom}, \texttt{Interval}, and \texttt{Calendar}. These components were then composed and aggregated to model complex interactions such as scheduling, notifications, and conflict handling.

\subsubsection{Design principles}\label{4hMyy7B7QRRagb14MwRDl}
\\
The design of the meeting scheduler reflects key SOLID principles:

\begin{itemize}
\item
\protect\phantomsection\label{W579bR1aZltjigsRR_V6n}
\textbf{Single Responsibility principle (SRP):} Each class is responsible for only one specific task (e.g., \texttt{Notification} only sends messages).
\item
\protect\phantomsection\label{vmLwVxOP11XIEg-ArukaJ}
\textbf{Open/Closed principle (OCP):} The system supports extending functionality (e.g., more notification types) without modifying existing code.
\item
\protect\phantomsection\label{-wvKqXzbIIxOWiFYnuWvV}
\textbf{Liskov Substitution principle (LSP):} All user types (organizers and participants) can be treated uniformly as \texttt{User} objects.
\item
\protect\phantomsection\label{qSossrMHfCvAmx599OQyH}
\textbf{Interface Segregation principle (ISP):} Interfaces are implied through class responsibilities without forcing unnecessary dependencies.
\item
\protect\phantomsection\label{yHHFfXtc5tMv3rYVsr52G}
\textbf{Dependency Inversion principle (DIP):} High-level modules like \texttt{MeetingScheduler} rely on abstractions like \texttt{Notification}, decoupling system parts.
\end{itemize}

\subsubsection{Design patterns}\label{QRqrLjuxrm0LHjMb3Wb6g}
\\
Several classical design patterns were applied to improve system maintainability:

\begin{itemize}
\item
\protect\phantomsection\label{L8zYBNEPviMs-4xOoehYz}
\textbf{Singleton pattern:} Ensures only one instance of \texttt{MeetingScheduler} manages the organization's meetings and rooms.
\item
\protect\phantomsection\label{T87cd-ebLSjFfEzqeChVp}
\textbf{Observer pattern:} Used for sending notifications to users when meetings are created, updated, or canceled, decoupling the messaging system from scheduling logic.
\end{itemize}

\subsection{Object interactions}\label{8Om6tod1k1sZKR-Gh7WcH}
\\
The following list describes how key objects interact dynamically during different system events:

\begin{itemize}
\item
\protect\phantomsection\label{AuYL5_C-CiAJyM-gz2AZM}
\textbf{Schedule meeting:}

\begin{itemize}
\item
The organizer provides meeting details.
\item
\texttt{MeetingScheduler} checks available \texttt{MeetingRoom} based on \texttt{Interval} and capacity.
\item
If available, the room is booked and a \texttt{Meeting} is created.
\item
\texttt{Meeting} is added to each participant's \texttt{Calendar}.
\item
Notifications are sent via \texttt{Notification} to all participants.
\end{itemize}
\item
\protect\phantomsection\label{l-GqCgGAoQdPphMx8TI2O}
\textbf{Cancel meeting:}

\begin{itemize}
\item
The organizer cancels the meeting.
\item
\texttt{MeetingScheduler} releases the \texttt{MeetingRoom}.
\item
\texttt{Meeting} is removed from participants' \texttt{Calendars}.
\item
Cancellation notifications are dispatched.
\end{itemize}
\item
\protect\phantomsection\label{6lzIXy7lMdpqycAbt_SFm}
\textbf{Respond to invite:}

\begin{itemize}
\item
Participant accepts or declines an invite.
\item
Their \texttt{Calendar} is updated accordingly.
\item
\texttt{Meeting} record is updated with their new \texttt{RSVPStatus}.
\end{itemize}
\item
\protect\phantomsection\label{RTDrj0VcO6x91fZ6rX29W}
\textbf{Update meeting:}

\begin{itemize}
\item
Organizer edits details (e.g., room, time).
\item
\texttt{MeetingScheduler} checks for conflicts, rebooks if needed.
\item
Calendars and notifications are updated.
\end{itemize}
\end{itemize}

\subsection{System workflow}\label{fYcYmXLzcf3TcnDEMNOu9}
\\
This section outlines the main workflows modeled in the system, derived from activity diagrams and use cases:

\begin{itemize}
\item
\protect\phantomsection\label{oltuSICAsnrxemLEc-iTQ}
\textbf{The organizer schedules a meeting:}

\begin{itemize}
\item
Opens calendar, selects time, sets agenda, and participants.
\item
\texttt{MeetingScheduler} checks and books an available room.
\item
Meeting is created, calendars updated, notifications sent.
\end{itemize}
\item
\protect\phantomsection\label{-LWUxXtDVtfg2sZUAUeZq}
\textbf{Participant responds to an invitation:}

\begin{itemize}
\item
Opens calendar, views meeting invite.
\item
Accepts or declines the meeting.
\item
Calendar and RSVP status are updated.
\end{itemize}
\item
\protect\phantomsection\label{Fh5YPi_eKZtqL_fT7bnpG}
\textbf{The organizer adds a participant:}

\begin{itemize}
\item
A new participant is added to the meeting.
\item
Their calendar is updated, and they receive an invitation.
\end{itemize}
\item
\protect\phantomsection\label{QCmKhba1sjW7EBEwz5Tz6}
\textbf{The organizer removes a participant:}

\begin{itemize}
\item
Participant is removed from the meeting.
\item
Calendar entry is deleted, and a cancellation notice is sent.
\end{itemize}
\item
\protect\phantomsection\label{3t2ybQuQ7utyI0b7MIHcL}
\textbf{The organizer cancels a meeting:}

\begin{itemize}
\item
The meeting is removed from all calendars.
\item
The room is released, and all participants are notified.
\end{itemize}
\end{itemize}

% End of section content